%==============================================================================
%  RAPPORT D'AVANCEMENT N°1 - Projet « Factures & Devis » - QUOVIL / NAFSIM Sarl
%  Compilation : pdflatex rapport_avancement_1.tex   (lancer 2 fois)
%  Moteur conseillé : pdfLaTeX (Overleaf, TeXstudio, MiKTeX, TeX Live)
%==============================================================================

%------------------------------------------------------------------------------
% 1. CLASSE DU DOCUMENT
%    article = document simple ; 11pt = taille du texte ; a4paper = format A4
%------------------------------------------------------------------------------
\documentclass[11pt,a4paper]{article}

%------------------------------------------------------------------------------
% 2. PACKAGES (extensions de LaTeX)
%------------------------------------------------------------------------------
\usepackage[utf8]{inputenc}      % lire les accents tapés au clavier (é, è, à...)
\usepackage[T1]{fontenc}         % bon encodage des caractères accentués
\usepackage{lmodern}             % police vectorielle propre
\usepackage[french]{babel}       % langue française (césures, « », dates...)
\usepackage[margin=2.2cm]{geometry} % marges de la page
\usepackage{graphicx}            % insérer des images (\includegraphics)
\usepackage{xcolor}              % couleurs
\usepackage{array}               % tableaux avancés (colonnes personnalisées)
\usepackage{longtable}           % tableaux pouvant continuer sur plusieurs pages
\usepackage{booktabs}            % (facultatif) beaux filets de tableaux
\usepackage{colortbl}            % colorer les cellules / lignes de tableaux
\usepackage{enumitem}            % listes personnalisées
\usepackage{fancyhdr}            % en-têtes et pieds de page
\usepackage{hyperref}            % liens cliquables (URL, sommaire)

%------------------------------------------------------------------------------
% 3. COULEURS ET PARAMÈTRES GÉNÉRAUX
%------------------------------------------------------------------------------
\definecolor{bleuensp}{RGB}{31,78,121}    % bleu des titres
\definecolor{bleuentete}{RGB}{46,116,181} % bleu des en-têtes de tableaux
\definecolor{orangeensp}{RGB}{237,125,49} % orange (rappel du logo ENSPY)

\hypersetup{
  colorlinks=true,
  linkcolor=bleuensp,
  urlcolor=bleuentete,
  pdftitle={Rapport d'avancement n°1 - Groupe G03},
  pdfauthor={Groupe G03 - 3 GTEL}
}

\setlength{\parindent}{0pt}               % pas de retrait en début de paragraphe
\setlength{\parskip}{6pt}                 % espace entre paragraphes
\renewcommand{\arraystretch}{1.25}        % lignes de tableaux un peu plus aérées

%------------------------------------------------------------------------------
% 4. COMMANDES PERSONNALISÉES
%------------------------------------------------------------------------------

% Colonnes de tableau : p{largeur} centrée (C) ou alignée à gauche (L)
\newcolumntype{C}[1]{>{\centering\arraybackslash}p{#1}}
\newcolumntype{L}[1]{>{\raggedright\arraybackslash}p{#1}}

% Cellule d'en-tête de tableau : texte blanc, gras, centré
\newcommand{\entete}[1]{\textcolor{white}{\textbf{#1}}}

% \imageouboite{fichier}{largeur}
% -> insère l'image si le fichier existe, sinon affiche une boîte à la place.
%    Ainsi le document COMPILE TOUJOURS, même sans les images.
\newcommand{\imageouboite}[2]{%
  \IfFileExists{#1}%
    {\includegraphics[width=#2]{#1}}%
    {\fbox{\parbox[c][3cm][c]{#2}{\centering\small\textit{Image manquante :}\\\texttt{\detokenize{#1}}}}}%
}

% Titre de section stylé (on redéfinit l'apparence de \section)
\usepackage{titlesec}
\titleformat{\section}{\Large\bfseries\color{bleuensp}}{\thesection-}{0.5em}{}[\titlerule]
\titleformat{\subsection}{\large\bfseries\color{bleuensp}}{\thesubsection}{0.5em}{}

% Numérotation des sections : 1, 2, 3...  ; sous-sections : a, b, c...
\renewcommand{\thesection}{\arabic{section}}
\renewcommand{\thesubsection}{\alph{subsection})}

%------------------------------------------------------------------------------
% 5. EN-TÊTE ET PIED DE PAGE (pages autres que la couverture)
%------------------------------------------------------------------------------
\pagestyle{fancy}
\fancyhf{}                                              % on vide tout
\fancyhead[L]{\small Rapport d'avancement n\textsuperscript{o}\,1 -- Groupe G03}
\fancyhead[R]{\small ENSPY -- 3 GTEL}
\fancyfoot[C]{\small Page \thepage}
\renewcommand{\headrulewidth}{0.4pt}

%==============================================================================
%                              DÉBUT DU DOCUMENT
%==============================================================================
\begin{document}

%------------------------------------------------------------------------------
% PAGE DE COUVERTURE
%------------------------------------------------------------------------------
\begin{titlepage}
\thispagestyle{empty}   % pas d'en-tête sur la couverture
\centering

\vspace*{0.5cm}
% Logo de l'ENSPY : placez le fichier logo_enspy.png à côté de ce .tex
\imageouboite{logo_enspy.png}{3.2cm}

\vspace{0.6cm}
{\color{bleuentete}\large RÉPUBLIQUE DU CAMEROUN -- \textit{Paix, Travail, Patrie}}

\vspace{0.4cm}
{\color{bleuentete}\large UNIVERSITÉ DE YAOUNDÉ 1}

\vspace{0.2cm}
{\large\textbf{ECOLE NATIONALE SUPERIEURE POLYTECHNIQUE\\DE YAOUNDÉ}}

\vspace{1.5cm}
{\normalsize\textbf{Cours de Base de Données -- Semestre 1 -- 2026 -- 2027}}

\vspace{0.3cm}
{\normalsize\textbf{3 GTEL}}

\vspace{1.5cm}
{\fontsize{30}{36}\selectfont\color{bleuensp} Rapport d'avancement n\textsuperscript{o}\,1\par}

\vspace{0.6cm}
{\large Projet « Factures \& Devis »}

\vspace{0.3cm}
{\large QUOVIL -- pour l'entreprise fictive « NAFSIM Sarl »}

\vspace{1.5cm}
{\normalsize Groupe : G03}

\vspace{0.3cm}
{\normalsize Période couverte : du 26/10/2026 au 09/11/2026}

\vspace{0.3cm}
{\normalsize Date de dépôt : 09/11/2026}

\vspace{0.3cm}
{\normalsize Dépôt Git : \url{https://github.com/AlphonseNoumbissi/groupe03.git}}

\vfill
\end{titlepage}

%------------------------------------------------------------------------------
% SECTION 1 : GROUPE & RÔLES
%------------------------------------------------------------------------------
\section{Groupe \& Rôles}

% longtable : tableau qui peut passer à la page suivante.
% Colonnes : N° | Noms & Prénoms | Matricule | Filière | Rôle | Taux
\begin{longtable}{|C{0.8cm}|L{4.6cm}|C{1.9cm}|C{1.7cm}|C{3.2cm}|C{1.5cm}|}
\hline
\rowcolor{bleuentete}
\entete{N\textsuperscript{o}} & \entete{Noms \& Prénoms} & \entete{Matricule} &
\entete{Filière} & \entete{Rôle} & \entete{Taux} \\
\hline
\endhead
1 & ALIMA KASSA Jean & / & GTEL & / & 0\,\% \\ \hline
2 & DJOUMA NGONGANG N.J. & 21p181 & GTEL & / & 0\,\% \\ \hline
3 & FONTSACK TASSADONG Hervé & 24p122 & GTEL & Développeur PDF / front & 23\,\% \\ \hline
4 & IMRAHNATOU YAYA & 26q375 & GTEL & Qualité \& documentation & 23\,\% \\ \hline
5 & MBILEU MBAKOP AMAYO Anaïs & 26q373 & GTEL & Concepteur BD / DBA & 24\,\% \\ \hline
6 & NOUMBISSI TEKEU Idriss Alphonse & 24p290 & GTEL & Chef de projet & 25\,\% \\ \hline
7 & SAAFOMAH TEZEMBON Zidane & 24p504 & GTEL & Développeur backend & 5\,\% \\ \hline
\end{longtable}

\textbf{Somme des taux = 100\,\%}

%------------------------------------------------------------------------------
% SECTION 2 : IDENTITÉ DE L'ENTREPRISE
%------------------------------------------------------------------------------
\section{Identité de l'entreprise}

NAFSIM SARL est une entreprise fictive spécialisée dans l'installation et la
maintenance des réseaux informatiques. Pour suivre son activité commerciale,
elle souhaite disposer d'une application permettant d'enregistrer les clients,
de préparer les devis, d'émettre les factures et de suivre les paiements.

\begin{longtable}{|C{0.8cm}|L{3.6cm}|L{10.2cm}|}
\hline
\rowcolor{bleuentete}
\entete{N\textsuperscript{o}} & \entete{Élément} & \entete{Exigences} \\
\hline
\endhead
1 & Raison sociale et sigle & NAFSIM SARL. \\ \hline
2 & Forme juridique et capital & Société à responsabilité limitée (SARL).
    Quatre millions huit cent mille FCFA. \\ \hline
3 & Logo & Voir figure~1 \\ \hline
4 & Cachet & Voir figure~2 \\ \hline
5 & Signature & Voir figure~3 \\ \hline
6 & Slogan / Activité & « Votre réseau, notre solution » \\ \hline
7 & Adresse & Quartier Melen, Yaoundé, BP 39393 Yaoundé \\ \hline
8 & Téléphones & +237 640203015 / +237 694768374 \\ \hline
9 & Email et site web & contact@nafsim.example \\ \hline
10 & Identifiants fiscaux & NIU : P098765432109 (fictif),
     RCCM : RC/YAO/2026/A/4587 (fictif) \\ \hline
11 & Coordonnées bancaires &
     \begin{itemize}[leftmargin=*,nosep]
       \item Banque : Afriland First Bank (compte 09642839642, nom du compte
             NAFSIM Sarl)
       \item Mobile Money : +237 698899899, nom NAFSIM Sarl
     \end{itemize} \\ \hline
12 & Conditions générales &
     \begin{itemize}[leftmargin=*,nosep]
       \item Validité des devis : 30 jours à compter de la date d'émission
       \item Délai de paiement : 75\,\% à la commande, solde à la livraison
       \item Mode de paiement : virement bancaire ou Mobile Money
       \item Pénalités de retard : en cas de retard, 10\,\% du montant total
             par semaine de retard ; à plus d'un mois, 50\,\% du montant total
             après 2 semaines
     \end{itemize} \\ \hline
\end{longtable}

%------------------------------------------------------------------------------
% FIGURES (logo, cachet, signature)
% Placez logo.png, cachet.png et signature.png dans le même dossier que ce
% fichier. S'ils n'existent pas, une boîte grise s'affiche (pas d'erreur).
%------------------------------------------------------------------------------
\begin{center}
\imageouboite{logo.png}{6cm}\\[2pt]
Fig1 : Logo

\vspace{0.6cm}
\imageouboite{cachet.png}{5cm}\\[2pt]
Fig2 : Cachet

\vspace{0.6cm}
\imageouboite{signature.png}{5cm}\\[2pt]
Fig3 : Signature
\end{center}

%------------------------------------------------------------------------------
% SECTION 3 : AUDIT DU CLASSEUR EXCEL
%------------------------------------------------------------------------------
\section{Audit du classeur Excel}

\begin{longtable}{|C{1cm}|L{7.4cm}|L{6.2cm}|}
\hline
\rowcolor{bleuentete}
\entete{N\textsuperscript{o}} & \entete{Anomalie constatée dans le document} &
\entete{Nom du problème (vocabulaire BD)} \\
\hline
\endhead
A1 & Les calculs intermédiaires et totaux (« 1\,060\,000 + 480\,000 =
     1\,540\,000 », remise, TVA, TTC) sont saisis sous forme de texte
     explicatif dans une cellule fusionnée au lieu d'être générés
     dynamiquement.
   & Données dérivées stockées/saisies au lieu d'être calculées (risque
     d'incohérence lors d'une modification). \\ \hline
A2 & La colonne Quantité de la section II contient des expressions
     mathématiques non atomiques (« 6*3 » et « 5*2 ») représentant à la fois
     le nombre de jours et le nombre de personnes.
   & Violation de la première forme normale (1NF) / non-atomicité : plusieurs
     attributs distincts (nombre de jours et nombre de personnes) sont
     fusionnés dans une seule colonne. \\ \hline
A3 & Mélange de deux concepts différents dans un même tableau : des articles
     matériels (section I) et des prestations de service (section II) qui
     n'ont pas les mêmes attributs (ex. : la main-d'œuvre dépend d'un effectif
     et d'une durée).
   & Mauvaise modélisation entité-association / absence de spécialisation
     (regroupement forcé d'entités hétérogènes). \\ \hline
A4 & Aucun code ou référence produit n'est attribué aux équipements (ex. :
     « routeur d'entreprise MultiWAN »). Seule une longue description
     textuelle est utilisée.
   & Absence d'identifiant (clé primaire) pour les articles et absence de
     référentiel (catalogue). \\ \hline
A5 & Les numéros de ligne « 1 » et « 2 » sont utilisés en double dans le même
     document (lignes 1 et 2 pour les fournitures, puis de nouveau lignes 1
     et 2 pour la main-d'œuvre).
   & Absence de clé primaire composite unique sur les lignes du devis
     (id devis). \\ \hline
A6 & La numérotation utilise des chiffres romains (I, II) pour les titres de
     sections, mélangés avec des chiffres arabes.
   & Absence de contrainte de domaine et de typage strict sur la colonne
     d'indexation. \\ \hline
A7 & Le client est simplement indiqué comme « client fictif », sans
     identifiant, adresse, contact ni date d'émission ou de validité du devis.
   & Absence de clé étrangère (id client) et de contraintes NOT NULL sur les
     métadonnées obligatoires du devis. \\ \hline
A8 & Les taux de remise globale (2\,\%) et de TVA (19,25\,\%) sont écrits en
     texte libre dans l'en-tête.
   & Redondance de données et absence d'attributs structurés pour les
     paramètres de facturation. \\ \hline
\end{longtable}

%------------------------------------------------------------------------------
% SECTION 4 : DICTIONNAIRE DES DONNÉES (BROUILLON)
%------------------------------------------------------------------------------
\section{Dictionnaire des données (brouillon)}

\subsection{Entités et attributs (brouillon)}

% \small réduit la taille du texte pour que le tableau tienne bien
{\small
\begin{longtable}{|L{3.2cm}|L{4.2cm}|L{7.8cm}|}
\hline
\rowcolor{bleuentete}
\entete{Entité} & \entete{Rôle} & \entete{Attributs} \\
\hline
\endhead
ENTREPRISE & Informations de l'entreprise émettrice &
  id\_entreprise (identifiant), raison\_sociale, forme, adresse, téléphone,
  email, numero\_identification, capital\_social, logo, cachet, signature,
  devise \\ \hline
UTILISATEUR & Personnes autorisées à utiliser l'application &
  id\_utilisateur (identifiant), nom, prénom, identifiant\_connexion,
  mot\_de\_passe\_hache, rôle, date\_creation, actif \\ \hline
CLIENT & Particulier ou entreprise cliente &
  id\_client (identifiant), code\_client, type\_client, nom, prénom,
  raison\_sociale, téléphone, email, numero\_identification, date\_creation,
  archive \\ \hline
ADRESSE\_CLIENT & Adresses associées aux clients &
  id\_adresse (identifiant), type\_adresse, adresse\_complete, ville,
  pays \\ \hline
FAMILLE\_ARTICLE & Classification des produits et prestations &
  id\_famille (identifiant), code\_famille, libelle, description \\ \hline
UNITE\_MESURE & Unités de quantification &
  id\_unite (identifiant), code\_unite, libelle \\ \hline
ARTICLE & Produits et prestations proposés &
  id\_article (identifiant), reference, designation, description,
  prix\_unitaire\_HT, actif \\ \hline
DOCUMENT & Informations communes aux devis, factures et avoirs &
  id\_document (identifiant), type\_document, numero\_document,
  date\_creation, date\_emission, date\_validite, statut, remise\_valeur,
  total\_HT, montant\_remise, taux\_TVA, montant\_TVA, total\_TTC,
  motif\_exoneration \\ \hline
SECTION\_DOCUMENT & Sections d'un document commercial &
  id\_section (identifiant), titre, ordre\_affichage, sous\_total\_HT \\ \hline
LIGNE\_DOCUMENT & Articles insérés dans une section &
  id\_ligne (identifiant), designation\_enregistree, quantite,
  prix\_unitaire\_HT\_applique, montant\_HT \\ \hline
MAIN\_D\_OEUVRE & Prestations évaluées selon le temps et les techniciens &
  id\_main\_oeuvre (identifiant), designation, nombre\_jours,
  nombre\_personnes, prix\_journalier, montant\_HT \\ \hline
FACTURE & Informations propres à une facture &
  id\_facture (identifiant), date\_echeance \\ \hline
REGLEMENT & Paiements effectués par les clients &
  id\_reglement (identifiant), date\_reglement, montant, mode\_paiement,
  reference\_paiement \\ \hline
AVOIR & Correction commerciale d'une facture &
  id\_avoir (identifiant), motif \\ \hline
JOURNAL\_AUDIT & Historique des opérations importantes &
  id\_journal (identifiant), action, type\_objet, id\_objet, date\_heure,
  details \\ \hline
\end{longtable}
}

\subsection{Associations et cardinalités}

Le tableau présente les associations principales du modèle proposé. Les
associations portant sur DOCUMENT doivent tenir compte de son type (DEVIS,
FACTURE ou AVOIR).

% enumerate = liste numérotée automatiquement (1., 2., 3., ...)
\begin{enumerate}[leftmargin=*,itemsep=4pt]
\item L'entité ENTREPRISE est liée à l'entité UTILISATEUR par l'association
  « EMPLOYER ». Cardinalités proposées : ENTREPRISE (1, N) -- UTILISATEUR
  (1, 1). Une entreprise emploie plusieurs utilisateurs ; chaque utilisateur
  appartient à une entreprise.
\item L'entité ENTREPRISE est liée à l'entité DOCUMENT par l'association
  « GÉRER ». Cardinalités proposées : ENTREPRISE (0, N) -- DOCUMENT (1, 1).
  Une entreprise peut émettre plusieurs documents ; chaque document appartient
  à une entreprise.
\item L'entité CLIENT est liée à l'entité ADRESSE\_CLIENT par l'association
  « POSSÉDER ». Cardinalités proposées : CLIENT (0, N) -- ADRESSE\_CLIENT
  (1, 1). Un client peut avoir plusieurs adresses ; une adresse appartient à
  un client.
\item L'entité CLIENT est liée à l'entité DOCUMENT par l'association
  « CONCERNER ». Cardinalités proposées : CLIENT (0, N) -- DOCUMENT (1, 1).
  Un client peut avoir plusieurs documents ; chaque document concerne un
  client.
\item L'entité UTILISATEUR est liée à l'entité DOCUMENT par l'association
  « CRÉER ». Cardinalités proposées : UTILISATEUR (0, N) -- DOCUMENT (1, 1).
  Un utilisateur peut créer plusieurs documents ; chaque document a un
  responsable.
\item L'entité FAMILLE\_ARTICLE est liée à l'entité ARTICLE par l'association
  « CLASSER ». Cardinalités proposées : FAMILLE\_ARTICLE (0, N) -- ARTICLE
  (1, 1). Une famille classe plusieurs articles ; chaque article appartient à
  une famille.
\item L'entité UNITE\_MESURE est liée à l'entité ARTICLE par l'association
  « MESURER ». Cardinalités proposées : UNITE\_MESURE (0, N) -- ARTICLE
  (1, 1). Une unité peut être utilisée par plusieurs articles ; chaque article
  a une unité.
\item L'entité DOCUMENT est liée à l'entité SECTION\_DOCUMENT par
  l'association « COMPRENDRE ». Cardinalités proposées : DOCUMENT (0, N) --
  SECTION\_DOCUMENT (1, 1). Un document peut contenir plusieurs sections ;
  chaque section appartient à un document.
\item L'entité SECTION\_DOCUMENT est liée à l'entité LIGNE\_DOCUMENT par
  l'association « CONTENIR ». Cardinalités proposées : SECTION\_DOCUMENT
  (0, N) -- LIGNE\_DOCUMENT (1, 1). Une section peut contenir plusieurs
  lignes ; chaque ligne appartient à une section.
\item L'entité ARTICLE est liée à l'entité LIGNE\_DOCUMENT par l'association
  « ÊTRE RÉFÉRENCÉ DANS ». Cardinalités proposées : ARTICLE (0, N) --
  LIGNE\_DOCUMENT (1, 1). Un article peut apparaître sur plusieurs lignes ;
  chaque ligne référence un article du catalogue.
\item L'entité SECTION\_DOCUMENT est liée à l'entité MAIN\_D\_OEUVRE par
  l'association « COMPORTER ». Cardinalités proposées : SECTION\_DOCUMENT
  (0, N) -- MAIN\_D\_OEUVRE (1, 1). Une section peut détailler plusieurs
  prestations de main-d'œuvre.
\item L'entité DOCUMENT est liée à l'entité FACTURE par l'association
  « CORRESPONDRE À ». Cardinalités proposées : DOCUMENT (0, 1) -- FACTURE
  (1, 1). Une facture spécialisée correspond à un document de type FACTURE.
\item L'entité FACTURE est liée à l'entité REGLEMENT par l'association
  « RECEVOIR ». Cardinalités proposées : FACTURE (0, N) -- REGLEMENT (1, 1).
  Une facture peut recevoir plusieurs règlements ; chaque règlement concerne
  une facture.
\item L'entité DEVIS/DOCUMENT est liée à l'entité FACTURE par l'association
  « ÊTRE À L'ORIGINE DE ». Cardinalités proposées : DEVIS (0, N) -- FACTURE
  (0, 1). Un devis peut ne produire aucune facture ou en produire plusieurs
  selon les règles retenues ; une facture peut être liée à un devis.
\item L'entité DOCUMENT/FACTURE est liée à l'entité AVOIR par l'association
  « ÊTRE CORRIGÉ PAR ». Cardinalités proposées : FACTURE (0, N) -- AVOIR
  (1, 1). Une facture peut être corrigée par plusieurs avoirs ; chaque avoir
  vise une facture d'origine.
\item L'entité DOCUMENT est liée à l'entité AVOIR par l'association
  « CORRESPONDRE À ». Cardinalités proposées : DOCUMENT (0, 1) -- AVOIR
  (1, 1). Un avoir spécialisé correspond à un document de type AVOIR.
\item L'entité UTILISATEUR est liée à l'entité JOURNAL\_AUDIT par
  l'association « EFFECTUER ». Cardinalités proposées : UTILISATEUR (0, N) --
  JOURNAL\_AUDIT (1, 1). Un utilisateur peut générer plusieurs entrées
  d'historique ; chaque entrée identifie son auteur.
\end{enumerate}

%------------------------------------------------------------------------------
% SECTION 5 : RÉSUMÉ EXÉCUTIF
%------------------------------------------------------------------------------
\section{Résumé exécutif}

Le projet en est à ses prémices. Nous avons déjà identifié le nom de
l'entreprise, le nom de l'application, les éléments d'identité constituant
l'entreprise, ainsi qu'un dictionnaire brouillon des différentes entités et
associations existantes dans notre entreprise. Néanmoins, nous avons rencontré
des problèmes au niveau du logo, de la signature et du cachet, ces éléments ne
figurant pas parmi les identifiants de l'entreprise, ainsi qu'au niveau de
l'audit du tableau Excel. Le principal risque est que le rapport présente des
manques d'informations et de données.

Nous progressons à présent dans la conception du MCD, de l'UML et du MLD en
respectant les différents éléments indiqués dans le cahier des charges. Le
principal risque est la possibilité de ne pas pouvoir exécuter le déclencheur
(\textit{trigger}).

%------------------------------------------------------------------------------
% SECTION 6 : OBJECTIFS ET ÉTAT D'AVANCEMENT
%------------------------------------------------------------------------------
\section{Objectifs de la période et état d'avancement}

\begin{longtable}{|C{1.5cm}|L{8cm}|C{1.5cm}|C{3cm}|}
\hline
\rowcolor{bleuentete}
\entete{Réf} & \entete{Objectifs / Exigences} & \entete{Prévu} & \entete{État} \\
\hline
\endhead
EF-01 & Paramétrer l'entreprise émettrice (12 éléments, téléversement du
        logo, du cachet et de la signature) & R1 & Terminé \\ \hline
OS1 & Créer l'identité complète de l'entreprise fictive & R1 & Terminé \\ \hline
OS2 & Analyser l'existant et rédiger le dictionnaire des données & R1 &
      Terminé \\ \hline
OS3 & Construire le MCD (MERISE) et les diagrammes UML & R2 &
      En cours (60\,\%) \\ \hline
EF-05 & Calcul en temps réel des montants en ligne & R1 & Terminé \\ \hline
\end{longtable}

%------------------------------------------------------------------------------
% SECTION 7 : CONTRIBUTION INDIVIDUELLE
%------------------------------------------------------------------------------
\section{Contribution individuelle}

\begin{longtable}{|L{3.8cm}|L{6.6cm}|C{1.5cm}|C{1.7cm}|}
\hline
\rowcolor{bleuentete}
\entete{Membres} & \entete{Tâches réalisées} & \entete{Heures} &
\entete{Commits} \\
\hline
\endhead
NOUMBISSI TEKEU I. A. & Planification, maquette de l'en-tête, rédaction du
  rapport, compte rendu de réunion, gestion du dépôt Git & 13 & 5 \\ \hline
IMRAHNATOU YAYA & Audit du classeur, calcul en temps réel des montants en
  lignes & 3 & 2 \\ \hline
MBILEU MBAKOP A. A. & Dictionnaire des données et identité partielle
  (9 éléments de l'identité de l'entreprise) & 10 & 2 \\ \hline
FONTSACK TASSADONG H. & Création du logo, du cachet et de la signature, et
  téléversement & 5 & 3 \\ \hline
SAAFOMAH TEZEMBON Z. & -- & 0 & 0 \\ \hline
\end{longtable}

\end{document}
